El objetivo general es analizar la importancia de la administración de la responsabilidad social y la ética en en el sector empresarial, destacando su papel en la gestión moderna. Se busca así, comprender que la conducta responsable de una organización no solo beneficia a la sociedad, sino también a la propia empresa.

Entre los objetivos específicos se pueden mencionar: explicar qué significa la responsabilidad social y la ética empresarial, identificar los principales principios que deben orientar la actuación de las empresas, describir algunas prácticas concretas que permiten aplicar estos principios en la vida diaria de la organización y destacar los beneficios que obtiene la empresa cuando incorpora estas ideas en su modelo de gestión. También se pretende mostrar que la responsabilidad social y la ética son herramientas útiles para prevenir problemas, fortalecer la confianza y mejorar la sostenibilidad empresarial.
